% Paper v1; same-vendor review PASS; two isolated cross-vendor referees (Claude Opus) PASS, 2026-10-02.
\documentclass[10pt]{article}
\usepackage[hmargin=1in,vmargin=0.9in]{geometry}
\usepackage{amsmath,amssymb,amsthm}
\usepackage{booktabs}
\usepackage{microtype}
\usepackage[hidelinks]{hyperref}
\hypersetup{pdftitle={An explicit second-order bound for Sidon sets},
  pdfauthor={Haoyu Chen}, pdfsubject={Finite Sidon sets; Erdos Problem 30}}
\numberwithin{equation}{section}
\newtheorem{theorem}{Theorem}[section]
\newtheorem{lemma}[theorem]{Lemma}
\theoremstyle{remark}
\newtheorem{remark}[theorem]{Remark}
\makeatletter
\renewcommand\section{\@startsection{section}{1}{\z@}{-2.6ex \@plus -.8ex \@minus -.2ex}{1.3ex \@plus .2ex}{\normalfont\large\bfseries}}
\renewcommand\subsection{\@startsection{subsection}{2}{\z@}{-2.2ex \@plus -.6ex \@minus -.2ex}{1ex \@plus .2ex}{\normalfont\normalsize\bfseries}}
\makeatother

\title{An explicit second-order bound for Sidon sets}
\author{Haoyu Chen\\ Independent researcher}
\date{October 2, 2026}

\begin{document}
\maketitle

\begin{abstract}
Let $F(N)$ be the largest size of a Sidon subset of $\{1,\ldots,N\}$,
with diagonal pair sums included. We prove
$F(N)\le\sqrt N+(2\sqrt2/3)N^{1/4}+1$ for every integer $N\ge120^4$.
The proof uses weighted Erd\H{o}s--Tur\'an energy with
$h(t)=2(1-t)\mathbf1_{[0,1]}(t)$ and an exact boundary correction of $2/3$;
an elementary contraction estimate makes the remainder explicit.
Reported coefficients run from Lindstr\"om's $1$ through Hou--Zhao's
$0.94349\ldots$ to unrefereed repository claims near $0.94324$ and
$0.94301$; the present coefficient is $0.942809\ldots$.
Erd\H{o}s's conjecture $F(N)=\sqrt N+O_\varepsilon(N^\varepsilon)$
for every $\varepsilon>0$ remains untouched.
\end{abstract}

\noindent\textit{Keywords:} Sidon sets; distinct differences; finite energy;
explicit bounds; boundary correction.

\section{Introduction}\label{sec:introduction}

A Sidon set of integers has distinct unordered two-term sums, including
sums in which the two summands agree. Write $F(N)$ for the largest
cardinality of such a set in $\{1,\ldots,N\}$. Translation gives the same
maximum on $\{0,\ldots,N-1\}$. The problem of bounding $F(N)$ goes back
to Erd\H{o}s and Tur\'an~\cite{ET41}, who proved
$F(N)\le\sqrt N+O(N^{1/4})$. Lindstr\"om~\cite{Lind69} gave the explicit
bound $F(N)<\sqrt N+N^{1/4}+1$, valid for every positive integer $N$.
Cilleruelo~\cite{Cil10} reduced its additive constant to $1/2$.
The reproductions and proofs in Balogh--F\"uredi--Roy~\cite{BFR23}
also record these two explicit statements.

Table~\ref{tab:history} gives a source-qualified progression of second-order
coefficients at the scale $\sqrt N+cN^{1/4}$. Where a source states a
uniform $O(1)$ remainder or a stronger form, the table indicates it;
no additional remainder assertion is assigned here to an announcement.
A coefficient with an unspecified $O(1)$ is not a printed numerical pair
of additive constant and onset. Publication years and preprint years are
listed separately when they differ. The table includes unrefereed
announcements and repository claims as such, and does not certify their
computational certificates.

\begin{table}[ht]
\centering\small
\begin{tabular}{@{}p{0.31\textwidth}p{0.19\textwidth}p{0.43\textwidth}@{}}
\toprule
Source & Coefficient $c$ & Form and status\\
\midrule
Erd\H{o}s--Tur\'an (1941)~\cite{ET41}
& unspecified & $O(N^{1/4})$; journal article\\
Lindstr\"om (1969)~\cite{Lind69}
& $1$ & additive $+1$, every $N\ge1$; journal article\\
Cilleruelo (2010)~\cite{Cil10}
& $1$ & additive $+1/2$, every $N\ge1$; journal article\\
Balogh--F\"uredi--Roy (2021/2023)~\cite{BFR23}
& $0.998$ & no additive term for all sufficiently large $N$; journal article\\
O'Bryant (2022/2024)~\cite{OB24}
& $0.99703$ & no additive term for all sufficiently large $N$; journal article\\
Carter--Hunter--O'Bryant (2023/2025)~\cite{CHO25}
& $0.98183$ & $O(1)$ remainder; journal article\\
AlphaEvolve collaboration (2025)~\cite{AE25}
& $0.97633$ & rounded safe coefficient; unrefereed announcement\\
Hou--Zhao v1 (July 2026)~\cite{HZ26}
& $0.9460052\ldots$ & $O(1)$ remainder; preprint\\
Hou--Zhao v3 (September 2026)~\cite{HZ26}
& $0.9434925907\ldots$ & $O(1)$ remainder; preprint\\
JoshMadeiros (August 2026)~\cite{Josh26}
& $\approx0.94324445$ & $c<0.94325$; $O(1)$; unrefereed forum claim\\
Wu, q1 (August 2026)~\cite{Wu26}
& $\approx0.94324253$ & $c<0.94325$; $O(1)$; unrefereed repository claim\\
Wu, q2 (August 2026)~\cite{Wu26}
& $\approx0.94300617$ & $c<0.94301$; $O(1)$; unrefereed repository claim\\
This note
& $2\sqrt2/3$ & additive $+1$ for every $N\ge120^4$; review status below\\
\bottomrule
\end{tabular}
\caption{Selected coefficient history. The initial date is the preprint or
announcement date when a later journal date is also displayed. Decimal
entries retain the scope and status of their cited sources.}
\label{tab:history}
\end{table}

Carter--Hunter--O'Bryant~\cite{CHO25} state a diameter bound
$k^2-1.96365k^{3/2}-O(k)$ for a $k$-element Sidon set, yielding the safe
cardinality coefficient $0.98183$ in the table. The 2025 collaboration of
Carter, Georgiev, G\'omez-Serrano, Hunter, O'Bryant, Tao and Wagner reported
a diameter coefficient $1.9526463099204112$; half that value explains the
rounded safe entry $0.97633$~\cite{AE25}. These comparisons concern
asymptotic coefficients: the difference between diameter $N-1$ and
interval parameter $N$ does not affect that coefficient, but must be
retained when asserting a specified additive constant.

The present note gives a complete proof with the numerical additive
constant and onset printed in the theorem. It uses the autocorrelation
of a decreasing linear weight. An exact correction on a half-line has
excess mass $1/3$; joining two copies gives the boundary term $2/3$.
A quantitative elementary contraction controls the tails. The auxiliary
measure on a finite interval is signed, which is permitted because the
energy comparison is an $L^2$ Cauchy--Schwarz inequality. No theorem about
renewal processes is needed.

Erd\H{o}s's conjecture asks for
$F(N)=\sqrt N+O_\varepsilon(N^\varepsilon)$ for every $\varepsilon>0$;
see the problem record~\cite{EP30}. The estimate proved here retains the
exponent $1/4$ and does not resolve that conjecture. The source comparison
above makes no claim that the displayed list is exhaustive or that the
method has a universal optimality barrier.

\section{Statement}\label{sec:statement}

A finite set of integers \(A\) is a \textbf{Sidon set} if the map

\[
\{(a,b)\in A^2:a\le b\}\longrightarrow\mathbb Z,
\qquad (a,b)\longmapsto a+b
\]

is injective. In particular, the pairs \((a,a)\) are included: equality
\(a+b=c+d\), with \(a\le b\) and \(c\le d\), requires \(a=c\) and \(b=d\).

\begin{theorem}\label{thm:main} For every integer \(N\ge207\,360\,000\) and every Sidon set
\(A\subseteq\{1,\ldots,N\}\),

\begin{equation}
\boxed{|A|\le\sqrt N+\frac{2\sqrt2}{3}N^{1/4}+1.}
\label{eq:s1-1}
\end{equation}
\end{theorem}

All powers and square roots in this document are positive real ones. In
particular, there is no requirement that \(N\) be a fourth power. The onset is
\(207\,360\,000=120^4\).

\begin{lemma}[equivalent difference convention]\label{lem:1} The Sidon condition is
equivalent to

\begin{equation}
a-b=c-d\ne0,\qquad a,b,c,d\in A
\quad\Longrightarrow\quad (a,b)=(c,d).
\label{eq:s1-2}
\end{equation}

Here the differences are indexed by ordered pairs, and zero differences are
excluded. Equivalently, each positive integer is the difference \(a-b\), with
\(a,b\in A\) and \(a>b\), at most once.

\end{lemma}

\begin{proof} Suppose first that pair sums, including diagonal sums, are unique.
The equality \(a-b=c-d\ne0\) gives \(a+d=c+b\). Equality of the two unordered
pairs with multiplicities gives either \(a=c,d=b\), as required, or
\(a=b,d=c\). The latter possibility contradicts the nonzero difference.

Conversely, suppose \eqref{eq:s1-2} holds, and let \(a+b=c+d\). If \(a=c\), then
\(b=d\). Otherwise \(a-c=d-b\ne0\), so \eqref{eq:s1-2}, applied to the ordered pairs
\((a,c)\) and \((d,b)\), gives \(a=d,c=b\). Thus the unordered pairs with
multiplicities agree. Restricting both pairs to increasing order proves the
stated Sidon condition, including when a pair is diagonal. Finally, a
negative difference becomes a positive one on reversing both ordered pairs;
hence uniqueness of positive differences is equivalent to \eqref{eq:s1-2}. \end{proof}

Translation preserves these conditions. We may consequently replace \(A\)
by \(A-1\subseteq\{0,\ldots,N-1\}\) when proving \eqref{eq:s1-1}.

\section{Kernel and energy}\label{sec:kernel}

Lebesgue measure on the real line is denoted by \(dt\). For an integrable
function \(v\) and a measure \(\mu\), write

\[
(v*\mu)(x)=\int v(x-t)\,d\mu(t)
\]

whenever the integral exists. A finite signed real measure means one with
finite total variation \(\|\mu\|_{\rm TV}=|\mu|(\mathbb R)\).

Set

\begin{equation}
h(t)=2(1-t)\mathbf1_{[0,1]}(t),\qquad
\widetilde h(t)=h(-t),\qquad f=h*\widetilde h.
\label{eq:s2-1}
\end{equation}

The specified value \(h(0)=2\) will be convenient in a pointwise identity.

\begin{lemma}[kernel and Cauchy inequality]\label{lem:2} The function \(f\) is even and

\begin{equation}
f(x)=
\begin{cases}
\displaystyle\frac43-2|x|+\frac23|x|^3,&|x|\le1,\\
0,&|x|>1.
\end{cases}
\label{eq:s2-2}
\end{equation}

It is nonnegative, nonincreasing on \([0,\infty)\), and satisfies

\begin{equation}
f(0)=\|h\|_2^2=\frac43,
\qquad \int_{\mathbb R}h(t)\,dt=1,
\qquad \int_{\mathbb R}f(t)\,dt=1.
\label{eq:s2-3}
\end{equation}

For finite signed real measures define

\[
E(\mu,\nu)=\iint f(x-y)\,d\mu(x)\,d\nu(y).
\]

Then

\begin{equation}
E(\mu,\nu)=\int_{\mathbb R}(h*\mu)(t)(h*\nu)(t)\,dt,
\qquad E(\mu,\mu)\ge0,
\label{eq:s2-4}
\end{equation}

and

\begin{equation}
|E(\mu,\nu)|^2\le E(\mu,\mu)E(\nu,\nu).
\label{eq:s2-5}
\end{equation}

\end{lemma}

\begin{proof} The autocorrelation is even by a change of variable. For
\(0\le x\le1\), its value is

\[
4\int_0^{1-x}(1-s-x)(1-s)\,ds
=\frac43-2x+\frac23x^3;
\]

for \(x>1\) the two supports have no overlap of positive length. This proves
\eqref{eq:s2-2}. Nonnegativity also follows directly from the integral of the two
nonnegative factors. On \([0,1]\) the derivative is \(-2+2x^2\le0\), and
outside this interval the function is zero. Integrating \(h\) and \(h^2\)
gives \eqref{eq:s2-3}; the mass of \(f\) is \((\int h)^2=1\) by Tonelli's theorem.

For any finite signed measure, Cauchy--Schwarz with respect to \(|\mu|\)
gives

\[
\int |h*\mu|^2\le
\|\mu\|_{\rm TV}
\int\!\int |h(t-x)|^2\,d|\mu|(x)\,dt
=\|\mu\|_{\rm TV}^2\|h\|_2^2.
\]

Thus \(h*\mu\in L^2\). Moreover,
\(\int |h(t-x)h(t-y)|\,dt\le\|h\|_2^2\), so the corresponding triple
integral against \(|\mu|\) and \(|\nu|\) is finite. Fubini's theorem is
therefore applicable. Its inner integral is \(f(x-y)\), which proves
\eqref{eq:s2-4}. Applying the usual Cauchy--Schwarz inequality in real \(L^2\) proves
\eqref{eq:s2-5}. This argument allows either measure to be signed. \end{proof}

\section{The half-line correction}\label{sec:halfline}

Let \(U\) be the probability measure with density
\(v_1=\mathbf1_{[0,1)}\), and let \(U^{*0}=\delta_0\). Define

\begin{equation}
g_+=\frac12\sum_{n=0}^{\infty}U^{*n}.
\label{eq:s3-1}
\end{equation}

The following construction proves everything needed about this measure;
no renewal theorem is used.

\begin{lemma}[local finiteness and exact potential]\label{lem:3} The measure \(g_+\) is
locally finite, supported on \([0,\infty)\), and

\begin{equation}
h*g_+=\mathbf1_{[0,\infty)},
\qquad (f*g_+)(x)=1\quad(x\ge0).
\label{eq:s3-2}
\end{equation}

\end{lemma}

\begin{proof} For \(n\ge1\), let \(v_n\) be the density of \(U^{*n}\).
Convolution gives representatives satisfying, for \(t\ge0\),

\begin{equation}
0\le v_n(t)\le\frac{t^{n-1}}{(n-1)!}.
\label{eq:s3-3}
\end{equation}

For \(n=1\) this is the bound \(v_1\le1\). If it holds at \(n\), then

\[
v_{n+1}(t)=\int_0^{\min(1,t)}v_n(t-s)\,ds
\le\int_0^t\frac{(t-s)^{n-1}}{(n-1)!}\,ds
=\frac{t^n}{n!}.
\]

For every fixed \(R\ge0\), it follows that
\(U^{*n}([0,R])\le R^n/n!\). The sum in \eqref{eq:s3-1} consequently has finite mass
on every bounded interval.

For a direct proof of the first identity in \eqref{eq:s3-2}, put
\(H=\mathbf1_{[0,\infty)}\). The definitions give the pointwise identity
\(h/2=H-U*H\). Thus each finite partial sum telescopes:

\begin{equation}
\frac12\sum_{n=0}^m h*U^{*n}
=H-U^{*(m+1)}*H.
\label{eq:s3-4}
\end{equation}

For \(t\ge0\), the last term on the right is between zero and
\(t^{m+1}/(m+1)!\), by \eqref{eq:s3-3}, and this tends to zero. For \(t<0\), both
sides vanish. The nonnegative partial sums on the left converge to
\(h*g_+\), proving the identity. This also fixes its value at \(t=0\).

Tonelli's theorem allows association of the nonnegative convolutions.
For \(x\ge0\),

\[
(f*g_+)(x)
=(\widetilde h*(h*g_+))(x)
=\int_{-1}^0\widetilde h(y)H(x-y)\,dy
=\int_{-1}^0\widetilde h(y)\,dy=1.
\]

All the integrals against \(g_+\) here are locally finite because their
kernels have compact support. \end{proof}

Write

\begin{equation}
g_+=\frac12\delta_0+\frac12u(t)\mathbf1_{[0,\infty)}(t)\,dt,
\qquad u(t)=\sum_{n=1}^{\infty}v_n(t)\quad(t\ge0).
\label{eq:s3-5}
\end{equation}

Use the right-continuous representative of \(u\). In particular,

\begin{equation}
u(0)=1,\quad u(t)=e^t\ (0\le t<1),\quad u(1)=e-1.
\label{eq:s3-6}
\end{equation}

The left limit at 1 is \(e\), and must not be substituted for \(u(1)\) in
the recurrence below. To justify these regularity claims, the series in
\eqref{eq:s3-5} converges uniformly on compact intervals by \eqref{eq:s3-3}. For \(n\ge2\),
the density \(v_n(t)=\int_{t-1}^t v_{n-1}(s)\,ds\), with densities set
to zero for \(s<0\), is continuous. Only \(v_1\) has a jump at 1. For
\(0\le t<1\), the induction proving \eqref{eq:s3-3} gives equality. At \(t=1\),
\(v_1(1)=0\) and \(v_n(1)=1/(n-1)!\) for \(n\ge2\). These observations
prove \eqref{eq:s3-6} and the stated representative convention.

\begin{lemma}[quantitative convergence of the density]\label{lem:4} For almost every
\(t\ge0\),

\begin{equation}
|u(t)-2|\le(e-1)(3/4)^{\lfloor t\rfloor}
<2(3/4)^{\lfloor t\rfloor}.
\label{eq:s3-7}
\end{equation}

\end{lemma}

\begin{proof} Summing the convolution recurrence for the densities yields
\(u=v_1+U*u\). Thus

\begin{equation}
u(t)=\int_{t-1}^t u(s)\,ds\qquad(t\ge1).
\label{eq:s3-8}
\end{equation}

This includes \(t=1\) with the right-continuous value in \eqref{eq:s3-6}. On
\((1,\infty)\), this identity makes \(u\) locally absolutely continuous,
and \(u'(t)=u(t)-u(t-1)\) almost everywhere. Fix an integer \(n\ge1\).
Solving this first-order equation on \([n,n+1]\) gives

\[
u(n+x)=e^xu(n)-\int_0^x e^{x-y}u(n-1+y)\,dy
\qquad(0\le x\le1).
\]

Since \(u(n)=\int_0^1u(n-1+y)\,dy\), this becomes

\begin{equation}
u(n+x)=\int_0^1K_x(y)u(n-1+y)\,dy,
\qquad
K_x(y)=e^x-\mathbf1_{\{y\le x\}}e^{x-y}.
\label{eq:s3-9}
\end{equation}

For \(0\le x,y\le1\), the kernel is nonnegative, and

\[
\int_0^1 K_x(y)\,dy=e^x-\int_0^x e^{x-y}\,dy=1.
\]

For \(y\in[1/2,1]\), it also satisfies \(K_x(y)\ge1/2\). Indeed, if
\(y>x\), its value is \(e^x\ge1\). If \(y\le x\), its value is
\(e^{x-y}(e^y-1)\ge e^{1/2}-1>1/2\), using the exponential series.

Let \([m_n,M_n]\) be the smallest closed interval containing the essential
range of \(u\) on \([n,n+1)\). Here essential extrema ignore sets of
Lebesgue measure zero. Initially \([m_0,M_0]=[1,e]\). Equation \eqref{eq:s3-9} first
implies that these closed intervals are nested. For the width estimate,
every \(K_x(y)\,dy\) contains the common measure

\[
w(y)\,dy=\frac12\mathbf1_{[1/2,1]}(y)\,dy
\]

of mass \(1/4\). Put
\(c_n=\int_0^1w(y)u(n-1+y)\,dy\). The remaining nonnegative measure
has mass \(3/4\), so every value in \eqref{eq:s3-9} lies between
\(c_n+(3/4)m_{n-1}\) and \(c_n+(3/4)M_{n-1}\). Consequently,

\[
[m_n,M_n]\subseteq[m_{n-1},M_{n-1}],
\qquad M_n-m_n\le\frac34(M_{n-1}-m_{n-1}).
\]

The nested closed intervals therefore meet in a single real number \(c\),
and

\begin{equation}
|u(t)-c|\le(e-1)(3/4)^{\lfloor t\rfloor}
\label{eq:s3-10}
\end{equation}

almost everywhere. This also provides a uniform essential bound and
essential uniform convergence to \(c\) as \(t\to\infty\).

For \(t>1\), the atom in \eqref{eq:s3-5} contributes nothing to \(h*g_+(t)\), so
Lemma~\ref{lem:3} gives

\[
1=\frac12\int_0^1 h(s)u(t-s)\,ds.
\]

Letting \(t\to\infty\) in this integral, \eqref{eq:s3-10} and \(\int h=1\) imply
\(1=c/2\). Thus \(c=2\). Finally, \(e<3\): in its series
\(e=2+\sum_{j\ge2}1/j!\), the tail is strictly smaller than
\(\sum_{j\ge2}2^{-(j-1)}=1\). This proves \eqref{eq:s3-7}. \end{proof}

\begin{lemma}[size, tail, and mass of the correction]\label{lem:5} Define the signed
measure supported on \([0,\infty)\)

\begin{equation}
q=g_+-\mathbf1_{[0,\infty)}(t)\,dt
=\frac12\delta_0+r(t)\mathbf1_{[0,\infty)}(t)\,dt,
\qquad r(t)=\frac12u(t)-1.
\label{eq:s3-11}
\end{equation}

It is a finite signed measure, and, with \(\alpha=\log(4/3)\),

\begin{equation}
|r(t)|\le(3/4)^{\lfloor t\rfloor}\quad\hbox{almost everywhere},
\qquad \|q\|_{\rm TV}\le\frac92,
\label{eq:s3-12}
\end{equation}

\begin{equation}
|q|((L,\infty))\le6e^{-\alpha L}\qquad(L\ge1),
\label{eq:s3-13}
\end{equation}

and

\begin{equation}
q(\mathbb R)=\frac13.
\label{eq:s3-14}
\end{equation}

\end{lemma}

\begin{proof} The first bound follows from Lemma~\ref{lem:4}. Integration on unit intervals
then gives

\[
\|q\|_{\rm TV}\le\frac12+\sum_{n=0}^\infty(3/4)^n=\frac92.
\]

For \(L\ge1\), the atom at zero is outside \((L,\infty)\), and

\[
|q|((L,\infty))
\le\sum_{n=\lfloor L\rfloor}^\infty(3/4)^n
=4(3/4)^{\lfloor L\rfloor}
\le\frac{16}{3}e^{-\alpha L}
\le6e^{-\alpha L}.
\]

It remains to compute the signed mass, not just its absolute bound. For
\(s>0\), Tonelli's theorem and the product rule for Laplace transforms of
nonnegative convolutions give

\begin{equation}
\begin{aligned}
\int e^{-st}\,dg_+(t)
&=\frac12\sum_{n=0}^\infty
\left(\int_0^1e^{-st}\,dt\right)^n\\
&=\frac1{2\left(1-(1-e^{-s})/s\right)}.
\end{aligned}
\label{eq:s3-15}
\end{equation}

The geometric series converges since \(0<\int_0^1e^{-st}\,dt<1\).
Taylor's formula for the exponential, as \(s\downarrow0\), gives

\[
\phi(s):=1-\frac{1-e^{-s}}s
=\frac s2-\frac{s^2}6+O(s^3).
\]

This use of \(O(s^3)\) concerns only a limit: for \(0<s\le1\), Taylor's
remainder after the cubic term of \(e^{-s}\) has absolute value at most
\(s^4/24\), which supplies the displayed bound directly. Subtracting the
Laplace transform \(1/s\) of half-line Lebesgue measure gives

\[
\int e^{-st}\,dq(t)
=\frac1{2\phi(s)}-\frac1s
=\frac{s-2\phi(s)}{2s\phi(s)}
\longrightarrow\frac13.
\]

Because \(q\) has finite total variation and is supported on the
nonnegative half-line, \(|e^{-st}|\le1\) there. Dominated convergence for
the finite measure \(|q|\) shows that the limit is \(q(\mathbb R)\), proving
\eqref{eq:s3-14}. \end{proof}

\subsection{Joining two half-lines}\label{sec:finite}

\begin{lemma}[finite-interval energy bound]\label{lem:6} Let \(L\ge1\), and let \(\mu\)
be any finite positive measure supported on the closed interval \([0,L]\).
If \(k=\mu(\mathbb R)\), then

\begin{equation}
E(\mu,\mu)\ge
\frac{k^2}{L+\frac23+200e^{-\alpha L}},
\qquad \alpha=\log(4/3).
\label{eq:s4-1}
\end{equation}

\end{lemma}

\begin{proof} Let \(q_L\) be the pushforward of \(q\) by the reflection
\(t\mapsto L-t\), and set

\begin{equation}
\nu_L=\mathbf1_{[0,L]}(t)\,dt+q+q_L.
\label{eq:s4-2}
\end{equation}

This is a finite signed real measure, with

\begin{equation}
\nu_L(\mathbb R)=L+\frac23.
\label{eq:s4-3}
\end{equation}

Let \(g_-^L\) denote the same reflection of \(g_+\). As locally finite
measures,

\begin{equation}
\nu_L=g_++g_-^L-dt.
\label{eq:s4-4}
\end{equation}

Indeed, the two half-line Lebesgue measures sum to full-line Lebesgue
measure plus Lebesgue measure on \([0,L]\); endpoints do not change those
Lebesgue measures. The atoms of \(q\) at 0 and of \(q_L\) at \(L\) remain
included in \eqref{eq:s4-2}. Since \(f\) is compactly supported, \eqref{eq:s4-4} may be
convolved locally with \(f\). Lemma~\ref{lem:3} and evenness of \(f\) give, for
\(0\le x\le L\),

\begin{equation}
V_L(x):=(f*\nu_L)(x)=1+1-\int f=1.
\label{eq:s4-5}
\end{equation}

This holds also at both endpoints. For every real \(x\), nonnegativity
and mass one of \(f\), together with Lemma~\ref{lem:5}, imply

\begin{equation}
|V_L(x)|\le1+\|f\|_\infty
(\|q\|_{\rm TV}+\|q_L\|_{\rm TV})
\le1+\frac43\cdot9=13.
\label{eq:s4-6}
\end{equation}

Outside \([0,L]\), the only measure in \eqref{eq:s4-2} on \((L,\infty)\) is the
tail of \(q\); on \(( -\infty,0)\) it is the reflection of the same tail.
In particular there are no excluded endpoint atoms. Thus \eqref{eq:s3-13} gives

\[
|\nu_L|(\mathbb R\setminus[0,L])\le12e^{-\alpha L}.
\]

By \eqref{eq:s4-3}, \eqref{eq:s4-5}, and \eqref{eq:s4-6},

\begin{equation}
\begin{aligned}
\left|E(\nu_L,\nu_L)-\left(L+\frac23\right)\right|
&=\left|\int_{\mathbb R\setminus[0,L]}(V_L(x)-1)\,d\nu_L(x)\right|\\
&\le14\cdot12e^{-\alpha L}
=168e^{-\alpha L}
\le200e^{-\alpha L}.
\end{aligned}
\label{eq:s4-7}
\end{equation}

Also \(E(\mu,\nu_L)=\int V_L\,d\mu=k\). Lemma~\ref{lem:2} now gives

\[
k^2\le E(\mu,\mu)E(\nu_L,\nu_L)
\le E(\mu,\mu)\left(L+\frac23+200e^{-\alpha L}\right).
\]

The denominator is positive and \(E(\mu,\mu)\ge0\), proving \eqref{eq:s4-1}.
There is no assumption that \(\nu_L\) is positive. \end{proof}

\section{The discrete inequality}\label{sec:discrete}

\begin{lemma}[Sidon energy upper bound]\label{lem:7} For every integer \(N\ge1\),
every Sidon set \(A\subseteq\{0,\ldots,N-1\}\), and every real
\(0<T\le N\), writing \(k=|A|\) gives

\begin{equation}
\boxed{
k^2\le
\left(\frac NT+\frac23+200e^{-\alpha N/T}\right)
\left(T+\frac43k\right).
}
\label{eq:s5-1}
\end{equation}

\end{lemma}

\begin{proof} Put \(L=N/T\ge1\) and
\(\mu=\sum_{a\in A}\delta_{a/T}\). Its mass is \(k\), and its support
lies in \([0,L]\). In fact its right endpoint is at most \((N-1)/T\);
using the larger interval of length \(N/T\) is permitted in Lemma~\ref{lem:6}.

Let \(D_+(A)\) denote the set of positive differences of elements of
\(A\). Lemma~\ref{lem:1} says that each such difference occurs for exactly one
ordered pair. The zero differences are precisely the \(k\) diagonal
ordered pairs. Evenness of \(f\) therefore gives the exact identity

\begin{equation}
E(\mu,\mu)=\frac43k+2\sum_{d\in D_+(A)}f(d/T)
\le\frac43k+2\sum_{d=1}^{\infty}f(d/T).
\label{eq:s5-2}
\end{equation}

The inequality uses nonnegativity of \(f\). For every integer \(d\ge1\),
monotonicity on the positive half-line implies

\[
f(d/T)\le T\int_{(d-1)/T}^{d/T}f(s)\,ds.
\]

Summing, and using evenness and \(\int f=1\), yields

\[
\sum_{d=1}^{\infty}f(d/T)
\le T\int_0^\infty f(s)\,ds=\frac T2.
\]

Only finitely many summands can be nonzero. It follows that
\(E(\mu,\mu)\le T+4k/3\). Combining this with \eqref{eq:s4-1} proves \eqref{eq:s5-1}.
This proof applies also to the empty set and to singleton sets. \end{proof}

\section{Explicit constants and onset}\label{sec:onset}

\begin{proof}[Proof of Theorem~\ref{thm:main}]

Let

\begin{equation}
x=N^{1/4},\quad T=\sqrt2\,x^3,\quad
\gamma=\frac{2\sqrt2}{3},\quad
\varepsilon=200e^{-\alpha x/\sqrt2}.
\label{eq:s6-1}
\end{equation}

For \(x\ge120\) the interval parameter \(L=N/T=x/\sqrt2\) is at least
one, so Lemma~\ref{lem:7} applies. Expanding its right side gives
\(P_\varepsilon(k)\le0\), where

\begin{equation}
P_\varepsilon(z)=
z^2-\left(\gamma x+\frac89+\frac43\varepsilon\right)z
-x^4-\gamma x^3-\sqrt2\,\varepsilon x^3.
\label{eq:s6-2}
\end{equation}

For clarity, the coefficient of \(z\) comes from
\((4/3)(x/\sqrt2+2/3+\varepsilon)\), and the constant term comes from
\(-(x/\sqrt2+2/3+\varepsilon)\sqrt2 x^3\).

We next establish a uniform bound for the error with all numerical
comparisons explicit. The integral formula for the logarithm gives

\[
\alpha=\int_1^{4/3}\frac{dt}{t}\ge\frac14.
\]

Since \(\sqrt2\le3/2\), it follows that
\(\alpha/\sqrt2\ge1/6\). The derivative of \(x e^{-x/6}\) is
\(e^{-x/6}(1-x/6)\), so it is nonpositive for \(x\ge6\). Consequently,
for every real \(x\ge120\),

\begin{equation}
\varepsilon x
\le200xe^{-x/6}
\le24000e^{-20}
<\frac{24000}{2^{20}}
=\frac{375}{16384}
<\frac1{32}.
\label{eq:s6-3}
\end{equation}

Here \(e>2\) follows from its series, and the last comparison is the
integer inequality \(375\cdot32=12000<16384\). Equivalently,
\(24000\cdot32=768000<1048576=2^{20}\).

This is the reason for the convenient choice \(x_0=120\), hence
\(N_0=120^4=207\,360\,000\): it makes the exponent \(x_0/6=20\) and
proves \eqref{eq:s6-3} by integer arithmetic. This onset is not claimed minimal.
The argument is uniform beyond the onset, rather than a check at one
point. More explicitly, the controlling envelope

\[
H(N)=200N^{1/4}\exp(-N^{1/4}/6)
\]

has derivative

\begin{equation}
H'(N)=50N^{-3/4}\exp(-N^{1/4}/6)
\left(1-\frac{N^{1/4}}6\right)\le0
\qquad(N\ge6^4).
\label{eq:s6-4}
\end{equation}

The actual quantity \(\varepsilon x\) is also decreasing there, since
\(\alpha/\sqrt2\ge1/6\). Thus no larger \(N\) can violate the error bound
used in \eqref{eq:s6-3}.

Put \(y=x^2+\gamma x+1\). Using \(\gamma^2=8/9\), direct expansion gives

\begin{equation}
P_0(y)=\frac{10}{9}x^2+\frac{\gamma}{9}x+\frac19\ge x^2.
\label{eq:s6-5}
\end{equation}

Also \(\gamma<1\), since \(\gamma^2=8/9<1\). For \(x\ge1\),
\(y\le x^2+x+1\le3x^2\). Hence \eqref{eq:s6-3} implies

\begin{equation}
\begin{aligned}
\varepsilon\left(\frac43y+\sqrt2\,x^3\right)
&<\frac1{32x}\left(4x^2+\frac32x^3\right)\\
&=\frac x8+\frac{3x^2}{64}
\le\frac{11x^2}{64}.
\end{aligned}
\label{eq:s6-6}
\end{equation}

The last step uses \(x\le x^2\), valid for \(x\ge1\). Combining
\eqref{eq:s6-2}, \eqref{eq:s6-5}, and \eqref{eq:s6-6} now proves, for every real \(x\ge120\),

\begin{equation}
P_\varepsilon(y)
=P_0(y)-\varepsilon\left(\frac43y+\sqrt2\,x^3\right)
>\frac{53}{64}x^2>0.
\label{eq:s6-7}
\end{equation}

The quadratic \(P_\varepsilon\) has leading coefficient 1 and strictly
negative constant term. Its two real roots have opposite signs: this also
follows directly from its positive discriminant and negative root product.
Therefore, on \([0,\infty)\), the condition \(P_\varepsilon(k)\le0\)
places \(k\) at or below the unique positive root. Since \(y>0\) and
\(P_\varepsilon(y)>0\), that root is strictly smaller than \(y\).
We have proved

\[
k<x^2+\gamma x+1
=\sqrt N+\frac{2\sqrt2}{3}N^{1/4}+1.
\]

Translating the original set as in Section~\ref{sec:statement} proves \eqref{eq:s1-1}. In particular,
the stated non-strict bound follows from a strict bound on the real right
side. \end{proof}

\section{Where \texorpdfstring{$2\sqrt2/3$}{2 sqrt(2)/3} comes from}\label{sec:constant}

Two exact constants determine this coefficient in the retained scalar
energy inequality:

\begin{equation}
a=f(0)=\frac43,
\qquad b=2q(\mathbb R)=\frac23.
\label{eq:s7-1}
\end{equation}

The first is the diagonal cost per point in \eqref{eq:s5-2}. The second is the sum
of the two half-line excess masses in \eqref{eq:s4-3}. Apart from the explicitly
controlled exponential error, \eqref{eq:s5-1} has right side

\begin{equation}
\left(\frac NT+b\right)(T+ak)
=N+bT+\frac{aNk}{T}+abk.
\label{eq:s7-2}
\end{equation}

To read its second-order scale, put \(N=x^4\),
\(T=t x^3\) for a fixed \(t>0\), and \(k=x^2+cx+O(1)\). The terms of
order \(x^3\) in the two sides of the inequality are respectively
\(2cx^3\) and \((bt+a/t)x^3\). The resulting upper coefficient is

\begin{equation}
c(t)=\frac12\left(bt+\frac at\right)
=\frac t3+\frac{2}{3t}.
\label{eq:s7-3}
\end{equation}

For \(t>0\),

\[
bt+\frac at-2\sqrt{ab}
=\left(\sqrt{bt}-\sqrt{a/t}\right)^2\ge0.
\]

Thus the minimum in \eqref{eq:s7-3} is
\(\sqrt{ab}=2\sqrt2/3\), attained precisely at
\(t=\sqrt{a/b}=\sqrt2\). This gives the scale chosen in \eqref{eq:s6-1}.
For fixed \(t\), the exponential term has argument \(-\alpha x/t\), so
it contributes no term at this order; Section~\ref{sec:onset} supplies its fully
effective treatment for the selected \(t\).

A precise limitation statement is also possible. It concerns \textbf{only the
scalar inequalities \eqref{eq:s5-1}, with this kernel and these retained energy
estimates}. It does not concern all arguments using this kernel, all
boundary corrections, or all Sidon methods.

To prove this limited statement, for each \(N>0\) let \(\kappa_N>0\) be
the solution of

\begin{equation}
\kappa_N=\sqrt N+\sqrt{ab\,\kappa_N}.
\label{eq:s7-4}
\end{equation}

Writing \(x=N^{1/4}\) and \(\gamma=\sqrt{ab}\) gives explicitly

\begin{equation}
\kappa_N
=\left(\frac{\gamma+\sqrt{\gamma^2+4x^2}}2\right)^2
=x^2+\gamma\sqrt{x^2+\gamma^2/4}+\frac{\gamma^2}{2}.
\label{eq:s7-5}
\end{equation}

For every \(T>0\), the elementary square inequality gives

\[
bT+\frac{aN\kappa_N}{T}\ge2\sqrt{abN\kappa_N}.
\]

Consequently,

\begin{equation}
\begin{aligned}
\left(\frac NT+b\right)(T+a\kappa_N)
&\ge N+2\sqrt{abN\kappa_N}+ab\kappa_N\\
&=\left(\sqrt N+\sqrt{ab\kappa_N}\right)^2
=\kappa_N^2.
\end{aligned}
\label{eq:s7-6}
\end{equation}

Equality occurs at \(T=\sqrt{aN\kappa_N/b}\). In particular,
\(\kappa_N\) satisfies the entire collection of these scalar
inequalities simultaneously, even if one allows every \(T>0\) and removes
the nonnegative exponential error. It therefore also satisfies \eqref{eq:s5-1} for
every allowed \(0<T\le N\).

Formula \eqref{eq:s7-5} gives the exact remainder identity

\begin{equation}
\kappa_N-x^2-\gamma x
=\frac{\gamma^2}{2}
+\frac{\gamma^3}{4\bigl(\sqrt{x^2+\gamma^2/4}+x\bigr)}.
\label{eq:s7-7}
\end{equation}

Thus \(\kappa_N=x^2+\gamma x+O(1)\). Even imposing integrality on a scalar
\(k\) does not change this leading limitation: for each fixed \(T\),
the inequality in \(k\) is a monic quadratic inequality with negative
constant term, so it holds on the full interval from zero to its positive
root. It therefore holds at \(\lfloor\kappa_N\rfloor\), simultaneously
for all \(T\), including with the exponential error restored. These
integers still equal \(x^2+\gamma x+O(1)\). For large \(N\) they also lie
between 0 and \(N\).

It follows that these scalar inequalities alone cannot imply
\(k\le\sqrt N+cN^{1/4}+C\) for any fixed \(c<\gamma\) and fixed \(C\):
the admissible scalar values \(\lfloor\kappa_N\rfloor\) eventually
violate that conclusion. These are scalar comparison values, \textbf{not
constructed Sidon sets}. Improving an earlier estimate, retaining more
information about the differences, or choosing another kernel lies outside
this limitation statement. No general method barrier is claimed.


\section{Remarks}\label{sec:remarks}

\subsection{A numerical kernel scan: non-rigorous evidence}
An exploratory discretization considered the minimum energy of signed
measures of mass one supported on $[0,L]$, for $L=4,8,16$, with
approximately $100$--$200$ grid points per unit interval. It used
$1/E_{\min}-L$ as an estimate of a boundary constant $b$ and compared
$\sqrt{ab}$, where $a=f(0)$. For the normalized weights
$h_p(t)=(p+1)(1-t)^p\mathbf1_{[0,1]}(t)$, the recorded approximate
scores were
\[
\begin{array}{c|rrrrrrr}
p&0&1/2&0.8&1&1.2&1.5&2\\\hline
\sqrt{ab}\text{ (reported)}&1.0000&0.9487&0.9435&0.9428&0.9433&0.9449&0.9487.
\end{array}
\]
The linear weight had a sampled boundary estimate $0.6687$, compared
with the exact $2/3$ proved above. The scan also reported scores near
$0.9470$ and $0.9446$ for normalized cosine and squared-cosine weights.
These are floating-point observations, without certified discretization
or rounding errors. They do not establish a local minimum in a
continuous family, optimality over kernels, or a limit of a general
method. None of these observations is used in the theorem or in the
scalar limitation statement of Section~\ref{sec:constant}.

\subsection{The same numerical constant in a different ruler problem}
Gupta and O'Bryant~\cite[Definition~2 and Theorem~3]{GO26} define an LM
ruler by requiring that a positive distance $d$ have at most $d-1$
representations. For an $n$-element LM ruler
$\{0=\ell_1<\cdots<\ell_n\}$ they prove
\[
\ell_n\ge\sqrt{8/9}\,(n-1)^{3/2}.
\]
Thus the number $\sqrt{8/9}=2\sqrt2/3$ occurs there as well. This is a
numerical coincidence between the stated constants: the multiplicity
condition, main scale, and conclusion differ. No implication between
that LM-ruler theorem and the Sidon cardinality bound of this note is
asserted.

\subsection{Analytic inputs and finite checks}
The proof uses only standard real analysis: Tonelli and Fubini under the
nonnegativity and absolute-integrability conditions verified in the text;
monotone and dominated convergence; Cauchy--Schwarz in real $L^2$;
absolute continuity of indefinite integrals of locally integrable
functions and its integrating-factor consequence; completeness of the
real numbers for nested compact intervals; and elementary exponential,
logarithmic, and quadratic calculations. The factorial majorant proves
compact uniform convergence of the density series. No external result
about Sidon sets, renewal processes, or extremal constants is an input
to the proof.

The accompanying audit script \texttt{check\_sidon\_bound.py} uses exact
fractions and integer square roots for rational enclosures and finite
sanity checks. A JavaScript translation was executed for the reported
finite checks, including maxima on intervals of lengths at most $60$;
the Python deliverable itself was not executed in the drafting environment.
These finite computations are not used to prove the theorem. In
particular, numerical grids do not replace the all-$N$ monotonicity
argument in Section~\ref{sec:onset}.

\section*{Provenance and verification status}
The argument was found by a clean-room GPT-6 Astra agent (OpenAI), then
expanded and examined in informed same-vendor adversarial reviews.
The source proof, the standalone expansion, and the finite-audit code
received same-vendor review; this is the completed mathematical review
status of the present draft. Cross-vendor referee reports by Claude
Opus (Anthropic): two isolated reports dated 2026-10-02, both \emph{PASS} with no repair required: referee~A re-derived every display and checked the exact potential of the capacity lemma numerically; referee~B reconstructed the main inequality independently from the statement and the constant, and tested every intermediate inequality on about $1.2\cdot10^5$ Sidon sets (all maximum Sidon sets for $N\le60$, Singer and Bose--Chowla sets for prime powers $q\le127$, Mian--Chowla prefixes, random and adversarial sets) without a violation. Both referees note that the onset $120^4$ is a convenient sufficient value: the same inequalities already give the stated bound for $N\ge4\,540\,589$ (numerical evaluation, not certified here).
No independent human referee report, Lean formalization, or other kernel
verification is claimed. The author takes responsibility for the
mathematics. Bibliographic statements and unrefereed numerical claims
are distinguished in the introduction; citing a claim does not certify
its proof or computational certificate.

\begin{thebibliography}{Josh26}\footnotesize\raggedright\setlength{\itemsep}{0.5pt}
\bibitem[ET41]{ET41} P.~Erd\H{o}s and P.~Tur\'an, On a problem of Sidon in additive number theory, and on some related problems,
\emph{J. London Math. Soc.} (1) \textbf{16} (1941), no.~4, 212--215,
\href{https://doi.org/10.1112/jlms/s1-16.4.212}{doi:10.1112/jlms/s1-16.4.212}.
\bibitem[Lind69]{Lind69} B.~Lindstr\"om, An inequality for $B_2$-sequences,
\emph{J. Combin. Theory} \textbf{6} (1969), no.~2, 211--212,
\href{https://doi.org/10.1016/S0021-9800(69)80124-9}{doi:10.1016/S0021-9800(69)80124-9}.
\bibitem[Cil10]{Cil10} J.~Cilleruelo, Sidon sets in $\mathbb N^d$,
\emph{J. Combin. Theory Ser. A} \textbf{117} (2010), no.~7, 857--871,
\href{https://doi.org/10.1016/j.jcta.2009.12.003}{doi:10.1016/j.jcta.2009.12.003}.
\bibitem[BFR23]{BFR23} J.~Balogh, Z.~F\"uredi and S.~Roy, An upper bound on the size of Sidon sets,
\emph{Amer. Math. Monthly} \textbf{130} (2023), no.~5, 437--445,
\href{https://doi.org/10.1080/00029890.2023.2176667}{doi:10.1080/00029890.2023.2176667};
preprint \href{https://arxiv.org/abs/2103.15850v2}{arXiv:2103.15850v2},
17 June 2021 (v1: 29 March 2021).
\bibitem[OB24]{OB24} K.~O'Bryant, On the size of finite Sidon sets,
\emph{Ukrains'kyi Matematychnyi Zhurnal} \textbf{76} (2024), no.~8, 1192--1206,
\href{https://doi.org/10.3842/umzh.v76i8.7858}{doi:10.3842/umzh.v76i8.7858};
preprint \href{https://arxiv.org/abs/2207.07800v2}{arXiv:2207.07800v2},
25 July 2022 (v1: 16 July 2022).
\bibitem[CHO25]{CHO25} D.~Carter, Z.~Hunter and K.~O'Bryant, On the diameter of finite Sidon sets,
\emph{Acta Math. Hungar.} \textbf{175} (2025), no.~1, 108--126,
\href{https://doi.org/10.1007/s10474-024-01499-8}{doi:10.1007/s10474-024-01499-8};
preprint \href{https://arxiv.org/abs/2310.20032v1}{arXiv:2310.20032v1},
30 October 2023.
\bibitem[AE25]{AE25} T.~Tao, Comment on \emph{Mathematical exploration and discovery at scale},
\emph{What's new}, 11 November 2025, 5:52 pm (site-displayed time),
\url{https://terrytao.wordpress.com/2025/11/05/mathematical-exploration-and-discovery-at-scale/#comment-689052}.
Unrefereed announcement; attribution and safe rounded coefficient also recorded at
\url{https://teorth.github.io/optimizationproblems/constants/5a.html}.
Accessed 2 October 2026.
\bibitem[HZ26]{HZ26} J.~Hou and H.~Zhao, An Improved Upper Bound for Finite Sidon Sets via Vector-Valued Smoothing,
preprint, \href{https://arxiv.org/abs/2607.01169v3}{arXiv:2607.01169v3},
4 September 2026 (v1: 1 July 2026; v2: 5 July 2026).
\bibitem[Josh26]{Josh26} J.~Madeiros (JoshMadeiros), An improved constant in the second-order term:
$h(N)\le N^{1/2}+0.94325N^{1/4}+O(1)$,
comment on Erd\H{o}s Problem \#30, 17 August 2026, 12:14 (site-displayed time),
\url{https://www.erdosproblems.com/forum/thread/30#post-8482}.
Unrefereed forum claim; accessed 2 October 2026.
\bibitem[Wu26]{Wu26} wustep (Stephen Wu), \emph{q2---grow free histograms from the q1 mix},
\texttt{maths} GitHub repository, 27 August 2026,
commit \texttt{da2440b68979b78d201182118d30ca418e3c2001},
\url{https://github.com/wustep/maths/blob/da2440b68979b78d201182118d30ca418e3c2001/problems/sidon-second-term/compute/q2/README.md}.
The q1 and q2 statements are also recorded in
\url{https://github.com/wustep/maths/blob/da2440b68979b78d201182118d30ca418e3c2001/problems/sidon-second-term/PROBLEM.md}.
Unrefereed repository claims; accessed 2 October 2026.
\bibitem[GO26]{GO26} A.~Gupta and K.~O'Bryant, Optimal Diameters of High Multiplicity $g$-Golomb Rulers,
preprint, \href{https://arxiv.org/abs/2605.14229v1}{arXiv:2605.14229v1},
14 May 2026.
\bibitem[EP30]{EP30} T.~F.~Bloom, Erd\H{o}s Problem \#30,
\url{https://www.erdosproblems.com/30}, accessed 2 October 2026.
\end{thebibliography}

\end{document}
